% Converted from 02-background.md.
% Required packages: \usepackage{amsmath}

\section{Background and Related Work}\label{background-and-related-work}

\subsection{Batch Auctions and the Closing Cross}\label{batch-auctions-and-the-closing-cross}

We adapt the Nasdaq Closing Cross {[}37{]} to clear one asset per batch on a grid of \(N\) price ticks. The same rule applies to each asset, whether an equity such as AAPL, a token such as ETH, or a futures contract.

\textbf{Order types.} We support buy and sell orders of three types, corresponding to limit-on-close, market-on-close, and imbalance-only orders. A \emph{limit order} specifies a limit price \(\ell\). At clearing price \(p^*\) a buy limit order is eligible if \(p^* \le \ell\), and a sell limit order if \(p^* \ge \ell\). A \emph{market order} specifies a quantity and is eligible at every price. An \emph{imbalance-only (IO) order} satisfies the same price condition as a limit order and may execute only on the short side. A buy IO can absorb excess supply, and a sell IO can meet excess demand.

\textbf{Clearing.} The first two phases select the price. The third determines how much can trade there.

\emph{Phase 1.} At each tick \(p\), the exchange sums eligible limit and market orders to obtain demand \(D(p)\) and supply \(S(p)\). The executable volume is \(V(p)=\min(D(p),S(p))\). Demand falls weakly with price and supply rises weakly, so the ticks attaining \(V_{max}=\max_p V(p)\) form an interval \([c,d]\), the \emph{plateau}. A positive maximum at a single tick determines the clearing price.

\emph{Phase 2.} If several ticks maximize a positive volume, the exchange minimizes \(\Delta(p)=|D(p)-S(p)|\) on \([c,d]\). A further tie is resolved by the at-price volume condition in Section 3.5.

\emph{Phase 3.} The exchange admits eligible IO volume on the short side up to the imbalance at \(p^*\). Better-priced eligible limit orders fill in full at the uniform price. If an imbalance remains, the exchange rations long-side volume at \(p^*\) pro rata.

A crossing batch yields the eight public values in the \emph{clearing receipt} (Section 3.1). A batch with zero executable volume at every tick yields a no-cross statement.

\subsection{Succinct Proofs}\label{succinct-proofs}

Zero-knowledge proofs let a verifier check a claim about private inputs {[}9, 24{]}. A succinct non-interactive argument of knowledge (SNARK) can have a logarithmic or constant-size proof in the size of the computation. Several SNARKs {[}17, 22, 32{]} use a PIOP {[}13{]}. The prover commits to polynomials and opens them at challenge points so the verifier can check identities among the evaluations. In KZG {[}29{]}, a commitment and an opening proof each occupy one group element. The universal reference string contains powers of a secret \(\tau\) in the exponent. A multiparty ceremony can generate it securely if at least one participant discards their secret contribution {[}39{]}. Fiat--Shamir {[}21{]} derives the verifier's challenges from a transcript that must bind the full public statement {[}6, 19{]}.

\textbf{Custom arguments.} A prover can execute a program on a verifiable machine, prove a compiled circuit, or work directly with polynomial constraints for the computation. Prior work considers the cost of this choice for ballot validity {[}28{]}. Zeeperio builds a PIOP for election tallies and reports substantial gains over circuit-based implementations {[}31{]}. The structure we exploit in an auction is the price grid. Each depth is a cumulative sum on that grid, and subsequent steps compare and select values from the same columns {[}40{]}.

\subsection{Related Work}\label{related-work}

\textbf{Secrecy-preserving auctions and exchanges.} Parkes et al. give a sealed-bid auction in which the auctioneer proves the result over encrypted bids {[}41{]}. Thorpe and Parkes prove properties of a private securities order book, including spread and depth {[}45{]}. In both constructions proof size grows with the number of orders. Other protocols distribute the computation among several parties. Bogetoft et al. use three-party secure computation to clear a double auction for Danish sugar-beet contracts {[}7{]}. Cartlidge, Smart, and Talibi Alaoui find that dark-pool volume matching is practical in their multiparty setting, while a continuous double auction is more costly {[}14{]}. FuturesMEX distributes a futures exchange among its traders {[}33{]}. Ngo et al. use witness key agreement for dark-pool trading {[}38{]}. A recent survey covers 22 protocol families {[}3{]}.

\textbf{Periodic and batch auctions.} Zhang and Ibikunle study how periodic auctions affect market quality in European equities {[}46{]}. Galal and Youssef construct a publicly verifiable periodic auction that keeps bids private by proving the order of committed values with aggregated Bulletproofs {[}12, 23{]}. FairTraDEX uses zero-knowledge set membership and an escrowed commit--reveal protocol to limit value extraction in a decentralized FBA. Orders become public when a batch ends {[}34{]}. We retain the exchange as operator and prove a closing-cross rule that includes IO orders and resolves ties across price ticks. For a fixed column layout, the clearing proof has a fixed number of commitments and openings. We know of no earlier construction with this combination of properties.

\textbf{Proofs by exchanges.} Provisions uses zero knowledge to prove solvency while protecting account balances {[}18{]}. Subsequent work identified weaknesses in several deployed solvency proofs {[}16{]}. Xiezhi uses a custom polynomial argument over KZG {[}20{]}. An exchange could use such a proof alongside ours to substantiate separate claims about its holdings and its clearing decisions.
